% secciones/01_introduccion/justificacion.tex
\subsection{Justificación}

Para los usuarios, el realizar una busqueda de peliculas puede resultar una tarea tediosa e incluso compleja cuando los usuarios no conocen el titulo de la pelicula o en general, no saben que pelicula quieren ver y simplemente necesitan recomendaciones de un genero en especifico. Para este tipo de consultas, es importante considerar tecnicas de procesamiento de lenguaje natural que permiten poder relacionar las necesidades expresadas en las preguntas del usuario con las caracteristicas de las peliculas en el catalogo.

Este proyecto se justifica porque permite comparar dos enfoques distintos para poder cumplir con el proposito del proyecto, TF-IDF y embeddings multilingues, ambos combinados con la similitud del coseno. Esta combinacion permite determinar si las diferencias en los modelos se reflejan en el desempeño del sistema. La eleccion de las tecnicas tambien se justifica ya que permite construir un prototipo funcional con herramientas de acceso libre y gratuito sin depender de servicios que requieran de una suscripcion o pago. 

Por otro lado, la incorporacion de un LLM permite transformar los resultados obtenidos en una respuesta mas completa y coherente para el usuario, esta capa se enfoca en la facilitacion de la presentacion de la informacion en pantalla, su funcionalidad se centra en unicamente redactar las respuestas a partir de la informacion obtenida en la busqueda, para de esta manera evitar el riesgo de que el modelo genere afirmaciones que no se encuentren respaldadas.

En general, el valor academico de este proyecto radica en que integra tecnicas de recuperacion de informacion y procesamiento de lenguaje, junto con sus metricas de evaluacion y el uso de Telegram como plataforma de interaccion con el usuario. La obtencion de resultados permite identificar las fortalezas y debilidades de cada enfoque sin asumir que una tecnica sera mejor que le otra.

De esta manera, al combinar todos estos factores se puede construir un prototipo funcional que permita a los usuarios realizar preguntas y recibir respuestas en español sobre peliculas.  